\section{定積分では上下端も換える}\label{節:置換と積分区間}

不定積分では最後に $x$ へ戻したが，定積分なら新しい変数のまま計算を終えられる。$t=g(x)$ と置いたときは，$x=a$ に対応する $t=g(a)$ を下端に，$x=b$ に対応する $t=g(b)$ を上端にする。大きい数を自動的に上端にするのではなく，元の端点との対応を保つ。

たとえば $t=1-x$ で $x$ が $0$ から $1$ へ動けば，$t$ は $1$ から $0$ へ動く。さらに $\dd t=-\dd x$ だから，
\[
 \int_0^1(1-x)^2\dd x
 =-\int_1^0 t^2\dd t
 =\int_0^1t^2\dd t=\frac13.
\]
微分から来た負号と，上下端を入れ換えたときの負号は，それぞれ別の理由で現れる。途中で片方だけを落とさないよう，対応と微分を先に書いておくとよい。

\begin{example}{途中で折り返す置換}
$\int_{-1}^{2}2x\cos(x^2)\dd x$ を求めよ。
\end{example}

$t=x^2$ と置けば $\dd t=2x\dd x$ であり，端点の対応は
\[
 \begin{array}{c|cc}
 x&-1&2\\ \hline
 t&1&4
 \end{array}
\]
である。したがって合成形の置換積分から
\[
 \int_{-1}^{2}2x\cos(x^2)\dd x
 =\int_1^4\cos t\dd t
 =\sin4-\sin1.
\]
角 $1,4$ はラジアンである。$x$ が $-1$ から $2$ まで増える間に，$t=x^2$ は一度 $1$ から $0$ へ減り，その後 $4$ へ増える。端点だけを見た式が正しいのは，この折り返しが定積分の向きに反映されているためである。

図\ref{図:置換の折返し}に対応を示す。$x=0$ で積分を分けて確かめると，
\[
\begin{aligned}
 \int_{-1}^{2}2x\cos(x^2)\dd x
 &=\int_1^0\cos t\dd t+\int_0^4\cos t\dd t\\
 &=-\int_0^1\cos t\dd t+\int_0^4\cos t\dd t
 =\int_1^4\cos t\dd t.
\end{aligned}
\]
往復した範囲の寄与が符号付きで打ち消し合っている。この計算は，曲線の下の面積を重複も含めて足す操作とは違う。

\begin{figure}[htbp]
 \centering
 \begin{tikzpicture}[x=1.6cm,y=0.8cm]
  \draw[-{Stealth}] (-1.35,0)--(2.35,0) node[right] {$x$};
  \draw[-{Stealth}] (0,-0.35)--(0,4.65) node[above] {$t$};
  \draw[thick,blue,domain=-1:2,samples=65] plot (\x,{\x*\x});
  \draw[dashed] (-1,0)--(-1,1)--(0,1);
  \draw[dashed] (2,0)--(2,4)--(0,4);
  \fill (-1,1) circle (1.5pt);
  \fill (2,4) circle (1.5pt);
  \node[below] at (-1,0) {$-1$};
  \node[below right] at (0,0) {$0$};
  \node[below] at (2,0) {$2$};
  \node[left] at (0,1) {$1$};
  \node[left] at (0,4) {$4$};
  \node[right] at (1.6,1.9) {$t=x^2$};
  \draw[-{Stealth},thick] (-0.88,0.85)--(-0.55,0.41);
  \draw[-{Stealth},thick] (1.12,1.38)--(1.43,2.17);
 \end{tikzpicture}
 \caption{$x$ は増え続けるが，$t=x^2$ は減少してから増加する。}
 \label{図:置換の折返し}
\end{figure}

